\section{задание: Рерайтинг и адаптация стиля}
\subsection{исходные данные}
\indent Тема исходника: \textbf{Новая функция} \\
\indent Целевая аудитория: \textbf{Пользователи} \\
\indent Требуемый стиль: \textbf{Простой} \\
\indent Ограничение: \textbf{ Твит (280 зн.) } \\
\indent Фокус: \textbf{Удобство} \\
\newline
\indent \textbf{Сгенерированный текст:} \\
\indent Представляем вашему вниманию релиз протокола взаимодействия с вычислительным ядром платформы. В текущий билд интегрирован модуль асинхронной обработки векторов состояний, который инициируется через вызов API-эндпоинта /api/v3/process с передачей массива байткодов. \\ 
\indent Ключевое улучшение заключается в имплементации механизма предиктивной выборки данных из L3-кэша, что позволило снизить задержки при исполнении тяжелых запросов. Внедренный алгоритм сжатия на лету (on-the-fly compression) на основе дельта-кодирования теперь автоматически применяется к полезной нагрузке (payload) для оптимизации пропускной способности канала. Логика обработки ошибок была расширена за счет добавления кодов состояния 429 Too Many Requests и 507 Insufficient Storage. \\
\indent Для обеспечения идемпотентности транзакций в заголовках запросов теперь обязательно требуется передавать UUID версии 4 в поле X-Request-ID. Система гарантирует консистентность данных в соответствии с требованиями ACID благодаря использованию двухфазного коммита (2PC) при взаимодействии с распределенным хранилищем. Для обратной совместимости со старыми клиентами оставлен "запасной" (fallback) эндпоинт, однако его поддержка будет прекращена в следующем мажорном релизе. Обратите внимание на обновленную схему аутентификации: теперь обязательна передача JWT-токена в каждом запросе, а не только при установлении сессии.
% \newpage
\subsection{базовый промпт}
\begin{verbatim}
					Перепиши это для {АУДИТОРИЯ}: {ТЕКСТ}
\end{verbatim}
предполагаю, что такого простого запроса должно хватить.
\subsection{журнал промпт-инжиниринга}
\noindent Промпт v1:
\begin{verbatim}
					Перепиши это для Пользователи: ...
\end{verbatim}
\newpage
\noindent Результат (YandexGPT):
\begin{figure}[H]
	\centering
	\includegraphics[width=17cm]{z2/a1.png}
\end{figure}
\newpage
\noindent Анализ ошибок: модель сгенерировала всё ещё технический документ. \\
Корректировка: дополнительная настройка запроса. \\
\newline
\noindent Промпт v2:
\begin{verbatim}
					Ты опытный айтишник.
					Тебя слушают: {АУДИТОРИЯ}.
					Перескажи текст,
					чтобы он был: {СТИЛЬ}, фокусируясь на: {ФОКУС}.
					Ограничение: {ОГРАНИЧЕНИЕ}.
					Текст: {ТЕКСТ}
\end{verbatim}
\noindent Результат (YandexGPT):
\begin{figure}[H]
	\centering
	\includegraphics[width=17cm]{z2/b1.png}
\end{figure}
\noindent Анализ ошибок: текст стал сильно короче, осталась только суть. \\
Корректировка: не требуется.
\newpage
\noindent Результат (DeepSeek):
\begin{figure}[H]
	\centering
	
\includegraphics[width=17cm]{z2/с1.png}
\end{figure}
\noindent Анализ ошибок: модель не использовала термины по типу UUID или JWT-токен \\
Корректировка: не требуется.